%!TEX root = ../main.tex
\section{Semantic Enrichment, Curation, and Routing}
\label{sec:curation}

After technical alignment, a pipeline must decide what behavior a record contains
and how it should influence training.  These are coupled problems: segmentation
determines the unit that can be labeled and valued, while a learned value signal
often reveals where a segment boundary or retry occurred.

\subsection{Why Episodes Are Too Coarse}

A teleoperated episode can include approach, grasp, transport, failed insertion,
retry, recovery, and terminal idle time.  One episode-level success bit maps all
transitions to the same target quality despite their different roles.  It also
prevents a long trajectory from being recomposed for a new instruction.  Modern
preparation therefore moves toward subtask segments and state transitions.

Boundary signals include gripper/contact changes, object motion, end-effector
velocity, change points in latent representations, language-aligned task progress,
and operator interventions.  Purely kinematic boundaries are cheap but can split
semantically atomic behavior or merge visually quiet planning with idle time.
Foundation-model labeling provides task semantics: NILS detects state changes and
uses pretrained models to produce language supervision without task-specific
training~\cite{blank2024scaling_2410_17772}.  Earlier, DIAL propagated a small
set of human trajectory descriptions through a vision--language model to augment
language supervision for a much larger robot dataset
~\cite{xiao2022robotic_2211_11736}.  Their strength is scalable relabeling;
their risk is a correlated error chain in object detection, state-change inference,
temporal grounding, and language generation.

A robust segmentation record should retain soft boundary confidence, the signals
that triggered it, and a hierarchy (episode $\supset$ phase $\supset$ primitive
$\supset$ transition).  Hard boundaries can then be sampled with tolerance during
training.  This representation supports three operations unavailable at episode
level: discard only an idle or corrupted interval; attach different language to
subtasks; and keep a failure prefix as context while routing a recovery suffix to
a recovery objective.

\subsection{Foundation-Model Labels Need Their Own Quality Control}

Visual-language models can infer objects, task descriptions, progress, and
success at a scale where exhaustive human annotation is impossible.  Their labels
are nevertheless proxy evidence, not ground truth.  Common failure modes include
hallucinating an occluded state, describing intention rather than achieved change,
using inconsistent granularity, missing brief contact, and inheriting geographic
or object-recognition bias from pretraining.

Preparation should therefore treat auto-labeling as selective prediction:
(i) generate structured fields before free text (objects, relations, before/after
state, action, outcome); (ii) query multiple views and temporal windows; (iii)
check cross-field constraints, such as an asserted placement requiring an object
relation change; (iv) calibrate confidence on a stratified human audit; and (v)
abstain or route disagreements to review.  Agreement among several prompts of the
same model is not independent evidence.  Better audits vary model families and
compare labels against measurable state or execution where available.

Label utility must also be tested downstream.  A fluent detailed description may
increase language entropy while fragmenting equivalent tasks into inconsistent
instructions.  Conversely, a coarse label can improve sharing but erase goal
parameters needed for precise control.  The appropriate granularity depends on
the training objective and stage; storing a label hierarchy preserves both.

\subsection{Failure Is a Type, Not a Quality Level}

Discarding all failed episodes removes unsafe final actions but also removes
near-success states, recovery examples, causal information about task boundaries,
and precisely the distribution encountered by an imperfect policy.  Keeping them
as expert demonstrations is equally problematic.  The resolution is
\emph{structure}: label failure cause and onset, separate valid prefixes and
recoveries, mark operator takeover and retry, and route each part to a compatible
objective.

SSDF uses relations between expert and imperfect data to select useful segments
from failures, showing in its own simulation and Franka experiments that imperfect
data need not be discarded wholesale~\cite{wu2024learning_2401_08957}.  ReTVL
uses retry structure to learn local value rather than assuming monotonic progress
~\cite{qin2026beyond_2606_24633}.  Demo-SCORE goes one step closer to the target
policy: a classifier trained on successful and unsuccessful online rollouts
filters unreliable demonstration strategies~\cite{chen2025curating_2503_03707}.
These methods resolve an apparent conflict in the literature.  ``Removing failed
data helps'' and ``failed data improves robustness'' can both hold because the
first often means removing mislabeled action targets, whereas the second uses
failure with outcome, boundary, retry, or recovery structure.

\subsection{Signals for Selection and Weighting}

\textbf{Rules and statistics.}  Range checks, idle ratio, path length, jerk,
smoothness, action saturation, success labels, and duplicate detection are
inexpensive and interpretable.  They are effective for known technical defects,
but a smooth trajectory may be confidently wrong and a corrective recovery may be
intentionally non-smooth.  Thresholds also correlate with episode length and task
difficulty.

\textbf{Representation and information.}  Embedding distance supports deduplication
and coverage; mutual information between state and action rewards demonstrations
that are both informative and predictable.  The mutual-information curation study
evaluates this signal on simulation and real robot data
~\cite{hejna2025robot_2502_08623}.  Crucially, its difficult Pen-in-Cup setting
slightly favors the full dataset over the filtered subset, consistent with lost
representation coverage.  This negative case is more informative than a universal
``less is more'' slogan: estimator quality, granularity, and retained support
matter.  Smoothness-driven RINSE similarly supplies a useful structural prior, but
its appropriate weight depends on whether a task contains impacts, retries, or
rapid corrections~\cite{kulkarni2026learning_2604_23000}.
Dataset-level entropy and learnability analyses broaden the diagnostic view
~\cite{xiao2025data_2511_09119}, while structure-aware selection such as SIEVE
operates over motion primitives and their transitions
~\cite{wu2026sieve_2607_06442}.  These signals are useful hypotheses about what
to inspect or retain, but neither a descriptive statistic nor a learned
structural score is policy utility until controlled retraining verifies it.

\textbf{Learned progress and redundancy.}  SCIZOR learns a task-progress signal
and deduplicates in joint state--action space, moving beyond image-only similarity
~\cite{zhang2025scizor_2505_22626}.  This is attractive at scale, but progress is
not monotonic in assembly, search, recovery, or tasks that temporarily undo a
state.  Progress-based filters require retry-aware calibration and should retain
low-confidence non-monotonic segments for targeted evaluation.

\textbf{Learner-aware value.}  Datamodel and influence approaches ask how a
training unit affects a particular learner or target set.  DataMIL estimates data
effects from subset-trained models~\cite{dass2025datamil_2505_09603}.  QoQ defines
quality through reduction of loss on validation demonstrations and aggregates
transition influence within trajectories~\cite{lee2026quality_2603_09056}.
ATHENA scales heterogeneous influence approximations to multi-task VLA curation
and evaluates both benchmark and real-robot tasks
~\cite{xu2026athena_2606_16208}.  These methods are closer to
Equation~\ref{eq:conditional-value}, but they exchange genericity for target
dependence: a narrow or biased validation set makes its preferences the definition
of quality; curvature approximations can be unstable; and repeated training,
rollouts, or influence computation can dominate saved training cost.

\begin{table*}[t]
\centering
\caption{Representative curation methods under a common schema. ``Real'' denotes a reported real-robot policy evaluation.}
\label{tab:curation-comparison}
\scriptsize
\setlength{\tabcolsep}{3.5pt}
\renewcommand{\arraystretch}{1.08}
\begin{tabularx}{\textwidth}{@{}p{2.25cm}p{2.05cm}p{2.8cm}p{1.6cm}ccX@{}}
\toprule
Method & Intervention & Guiding signal & Granularity & Human labels & Real & Principal caveat \\
\midrule
Data Quality~\cite{belkhale2023data_2306_02437} & Diagnose / construct & Action divergence and transition diversity & Dataset / transition & No & No & Formal lens, not a universal rank \\
SSDF~\cite{wu2024learning_2401_08957} & Recover failure segments & Expert--imperfect relation & Segment & No & Yes & Needs expert support \\
Re-Mix~\cite{hejna2024re_2408_14037} & Learn source weights & Weakest-group validation loss & Source / domain & Target split & Yes & Costly fits, abnormal-source failure \\
Mutual information~\cite{hejna2025robot_2502_08623} & Filter demonstrations & State--action mutual information & Episode & No & Yes & Omits causality, coverage trade-off \\
SCIZOR~\cite{zhang2025scizor_2505_22626} & Filter / deduplicate & Progress and state--action similarity & Transition / segment & No & Yes & Monotonic-progress assumption \\
DataMIL~\cite{dass2025datamil_2505_09603} & Select subsets & Response to subset training & Episode / task & No & Yes & Repeated training, learner dependence \\
Demo-SCORE~\cite{chen2025curating_2503_03707} & Filter strategies & Classifier from rollout outcomes & Episode / strategy & Outcome & Yes & Needs representative rollouts \\
QoQ~\cite{lee2026quality_2603_09056} & Influence selection & Validation--demonstration loss & Transition $\to$ episode & Validation set & Yes & Inherits target-set bias \\
ATHENA~\cite{xu2026athena_2606_16208} & Multitask selection & Heterogeneous influence & Sample / task & Target set & Yes & Approximation and target construction \\
ReTVL~\cite{qin2026beyond_2606_24633} & Retry-aware weighting & Retry-supervised value & Transition / segment & Retry event & Yes & Retry is not synonymous with low value \\
\bottomrule
\end{tabularx}
\end{table*}


\subsection{Detection Quality Is Not Policy Utility}

The controlled audit ``What Demonstration Curation Metrics Do to Your Policy''
explicitly separates detection of injected defects from the performance of a
policy trained after curation~\cite{bedi2026what_2606_10229}.  It finds that a
metric's ability to recognize an intended structural defect can disagree with
downstream utility and exposes episode-length confounding.  Its small task and
rollout scale limit generality, but the methodological warning is fundamental.
A detector can achieve high AUROC by identifying a visible artifact while
removing rare states needed for recovery; policy utility can improve even with
modest defect recall if the removed examples exert disproportionate gradients.

Accordingly, every curation paper should answer two separate questions:
\begin{enumerate}[leftmargin=*]
    \item Did the operation change the intended data property at a stated
    threshold, and what other properties changed with it?
    \item Under equal data access and compute, did that change improve a trained
    policy, on which target distribution, with what uncertainty?
\end{enumerate}
The first diagnoses a mechanism; only the second supports a utility claim.

\subsection{Mixtures, Sampling, and Role Routing}

Cross-source training adds a compositional decision: proportions affect both
gradient scale and how often rare capabilities are seen.  Uniform sampling
overweights large sources; uniform-by-dataset can over-amplify tiny noisy sources;
fixed hand weights become stale as the learner changes.  Re-Mix formulates
mixture-weight learning through distributionally robust optimization and reports
large gains over uniform or human mixtures in its own two-embodiment evaluation
~\cite{hejna2024re_2408_14037}.  The paper also identifies a revealing failure:
the learned objective can increase weight on abnormal action distributions.  The
optimization target therefore requires constraints or audits, not blind trust.

Language-model preparation offers transferable mechanisms, not transferable
quality definitions.  DoReMi learns domain weights from excess loss, while DSIR
resamples according to an importance estimate toward a target distribution
~\cite{xie2023doremi_2305_10429,xie2023data_2302_03169}.  In robotics, analogous
mixture updates must account for temporal correlation, embodiment, execution
cost, and physical validity.  A useful hierarchy is source $\rightarrow$ task
$\rightarrow$ episode $\rightarrow$ segment, with floors that protect rare
embodiments and safety-critical failures from collapse.

Routing is stronger than weighting.  Rather than asking whether an example is
globally good, route it to a role: broad but weakly aligned video to representation
pretraining; aligned robot action to behavior cloning; unsuccessful attempts with
outcome labels to value or recovery learning; and uncertain synthetic episodes to
verifier training.  Stage-conditioned routing explains why industrial recipes
often use a broad pretraining mixture followed by narrower embodiment and task
post-training, and why one fixed mixture is rarely optimal for the entire
lifecycle.

\subsection{Synthesis: The Precision--Coverage--Cost Frontier}

Across curation methods, stricter thresholds typically improve one observable
property while reducing coverage and data volume.  More diversity can initially
expand support but later inject incompatible conventions; smoothness can reject
noise but also reject corrective behavior; a target-aware estimator can improve
near-term loss but narrow long-term transfer.  The governing object is therefore
a Pareto frontier among correctness, coverage/diversity, target alignment, and
cost.  A robust engine should expose this frontier, calibrate cheap proxies with
periodic fixed-budget retraining, and preserve discarded data and reason codes so
decisions can be reversed when the learner or target changes.
